\documentclass[a4paper,12pt]{ctexart}

% Using ctexart with UTF-8 encoding
\usepackage{booktabs}
\usepackage{longtable}
\usepackage{array}
\usepackage{colortbl}
\usepackage{geometry}
\usepackage{xcolor}
\usepackage{fontspec}
\usepackage{indentfirst}
\usepackage{setspace}
\usepackage{amsmath} % Add amsmath for \text command
\usepackage{ulem} % For underlining text
\usepackage{mdframed} % For creating frames around content

% Page setup with narrower margins
\geometry{
  a4paper,
  margin=2cm,
  includehead
}

% Set fonts
\setCJKmainfont{Kaiti TC}  % Use Kaiti TC for main Chinese text
\setCJKsansfont{Kaiti TC}  % Use Kaiti TC for sans-serif Chinese text
\newCJKfontfamily\hanwangfont{HanWangKaiMediumChuIn}

% Define commands for vocabulary words with and without highlighting
\newcommand{\vocab}[1]{{\colorbox{yellow}{\hanwangfont #1}}}
\newcommand{\vocabnohi}[1]{{\fontsize{14pt}{16pt}\selectfont\hanwangfont #1}}

% Define a command for English translations/explanations
\newcommand{\eng}[1]{{\small\itshape #1}}

\begin{document}

\begin{center}
{\LARGE\bfseries 第四課 求職與面試}

\vspace{0.3cm}
{\large 學習目標（Tujuan Pembelajaran）}
\end{center}

\begin{enumerate}
\item 能夠討論自己想申請的工作類型與條件。Mampu mendiskusikan jenis pekerjaan dan persyaratan yang ingin Anda lamar.
\item 能夠正式、專業地回答面試相關問題。Mampu menjawab pertanyaan terkait wawancara secara formal dan profesional.
\end{enumerate}

\vspace{0.5cm}
\setstretch{1.3}

\begin{mdframed}[linewidth=1pt]
\noindent\vocab{面試官}：你好，歡迎你來參加\vocab{面試}，請你先\vocab{簡單}地介紹一下自己。

\noindent 林明華：您好，我叫林明華，\vocab{雙}木林，明天的明，中華的華。我是美國人，去年從成功大學企管系\vocab{畢業}。\vocab{畢業}以後，在一家\vocab{科技}公司\vocab{擔任}\vocab{業務助理}，上個月剛\vocab{離職}。

\noindent\vocab{面試官}：請問你為什麼\vocab{離職}？有什麼問題嗎？

\noindent 林明華：沒有什麼問題。\uline{我\vocab{離職}是因為想找一個更有\vocab{挑戰性}的工作}。

\noindent\vocab{面試官}：我們公司有很多外國\vocab{客戶}，你\vocab{認為}你的語言能力可以\vocab{勝任}嗎？

\noindent 林明華：我的語言能力不錯，\uline{除了英文、中文很流利之外，我還會說日文和一點西班牙文}。

\noindent\vocab{面試官}：你會說這麼多語言，真讓人\vocab{印象深刻}。可以請你\vocab{談談}你認爲自己做得最成功的一件事嗎？

\noindent 林明華：好的。我認爲自己從美國來台灣\vocab{留學}，\vocab{克服}了語言、文化的障礙，在台灣累積了許多經驗，而且在台灣交了很多好朋友，這是我做得最成功的一件事。

\noindent\vocab{面試官}：你\vocab{應徵}的\vocab{職位}是\vocab{行銷}\vocab{專員}，那你對我們公司的商品了解嗎？

\noindent 林明華：我知道\vocab{貴公司}\vocab{生產}手機和\vocab{平板電腦}。我對這\vocab{方面}很有興趣，常常注意\vocab{相關}新聞。

\noindent\vocab{面試官}：好的。你有沒有什麼想問我們的問題？

\noindent 林明華：我想更了解貴公司的上下班\vocab{制度}與\vocab{員工}\vocab{福利}，不知道能不能請您簡單介紹一下？

\noindent\vocab{面試官}：\uline{在上下班方面，時間比較\vocab{彈性}，員工也可以申請\vocab{遠端}工作}。至於福利，\uline{除了\vocab{健保}和\vocab{勞保}，我們公司每年還提供員工健康檢查的假期與\vocab{補助}}。另外，也會根據員工表現給\vocab{年終獎金}。

\noindent 林明華：我了解了，感謝您的回答。

\noindent\vocab{面試官}：好的，今天的面試就到這裡，我們會在一週內\vocab{通知}你面試\vocab{結果}。

\noindent 林明華：謝謝，希望有機會到貴公司服務。

\noindent\vocab{面試官}：謝謝，再見。
\end{mdframed}

\vspace{1cm}
\begin{center}
{\Large\bfseries 生詞（Kosakata）}
\end{center}

\begin{longtable}{lp{12cm}}
\vocabnohi{面試官} & n. pewawancara \\
\vocabnohi{簡單} & adj. sederhana, simpel \\
\vocabnohi{雙} & adj. dua, ganda (lawan dari 單) \\
\vocabnohi{畢業} & v. lulus, wisuda \\
\vocabnohi{科技} & n. teknologi \\
\vocabnohi{擔任} & v. menjabat sebagai \\
\vocabnohi{業務助理} & n. asisten bisnis \\
\vocabnohi{離職} & v. mengundurkan diri, berhenti kerja \\
\vocabnohi{挑戰性} & n. menantang \\
\vocabnohi{客戶} & n. klien, pelanggan \\
\vocabnohi{認為} & v. berpikir, berpendapat bahwa \\
\vocabnohi{勝任} & v. memenuhi syarat untuk \\
\vocabnohi{印象深刻} & terkesan \\
\vocabnohi{談} & v. berbicara \\
\vocabnohi{留學} & v. belajar di luar negeri \\
\vocabnohi{克服} & v. mengatasi \\
\vocabnohi{困難} & n. kesulitan \\
\vocabnohi{應徵} & v. melamar \\
\vocabnohi{職位} & n. posisi, jabatan \\
\vocabnohi{行銷} & n. pemasaran \\
\vocabnohi{專員} & n. komisioner \\
\vocabnohi{貴公司} & n. perusahaan Anda (lebih sopan) \\
\vocabnohi{生產} & v. memproduksi \\
\vocabnohi{平板電腦} & n. tablet \\
\vocabnohi{方面} & n. aspek \\
\vocabnohi{相關} & n. terkait, berhubungan \\
\vocabnohi{制度} & n. sistem \\
\vocabnohi{員工} & n. karyawan \\
\vocabnohi{福利} & adj. kesejahteraan \\
\vocabnohi{彈性} & adj. fleksibilitas \\
\vocabnohi{遠端} & n. jarak jauh \\
\vocabnohi{健保} & n. asuransi kesehatan (singkatan dari \vocabnohi{健康保險}) \\
\vocabnohi{勞保} & n. asuransi tenaga kerja (singkatan dari \vocabnohi{勞工保險}) \\
\vocabnohi{補助} & v. subsidi \\
\vocabnohi{年終獎金} & v. bonus akhir tahun \\
\vocabnohi{通知} & n. memberi tahu \\
\vocabnohi{結果} & n. hasil \\
\end{longtable}

\vspace{0.5cm}
\begin{center}
{\Large\bfseries 語法點（Catatan Tata Bahasa）}
\end{center}

\setstretch{1.2}
\begin{enumerate}
\item \textbf{我叫林明華，雙木林，明天的明，中華的華。}\\
Karena banyak karakter Mandarin memiliki pengucapan yang sama, penutur asli sering menambahkan informasi tambahan saat memperkenalkan nama mereka untuk menghindari kebingungan.\\
Salah satu cara umum adalah dengan mendeskripsikan struktur karakter, seperti "\vocab{雙木林}" (林 yang terdiri dari dua 木).\\
Cara lain adalah dengan menghubungkan karakter ke kata atau frasa yang familiar, seperti "\vocab{明天的明}" (明 dalam 明天).\\
Ini membantu pendengar memahami karakter mana yang dimaksud dengan cepat.

\item \textbf{我離職是因為想找一個更有挑戰性的工作。}\\
\text{[Hasil]} ……\vocab{是因為} \text{[Alasan]}…… \\
a. 我搬去台北是因為那裡的工作機會比較多，薪水也比較高。\\
b. 你做得不好，不是因為沒有能力，是因為不夠認真。

\item \textbf{除了英文、中文很流利之外，我還會說日文和一點西班牙文。}\\
Bagian antara \vocab{除了} dan \vocab{之外} adalah informasi yang sudah diketahui (informasi lama) dan bagian setelah \vocab{還} adalah informasi baru. \vocab{除了} dapat diikuti oleh kata benda, kata kerja, kata sifat atau frasa. \vocab{之外} dapat dihilangkan.\\
a. 除了健保和勞保，我們公司每年還提供員工健康檢查的假期與補助。\\
b. 我媽媽除了中國菜之外，還會做法國菜和泰國菜。

\item \textbf{這類工作年終獎金比較高，在上下班方面也比較彈性。}\\
(在）……\vocab{方面} berkenaan dengan, mengenai\\
\vocab{方面} berarti bidang, lingkup atau kategori.（在）……\vocab{方面} digunakan ketika pembicara ingin mengungkapkan pendapat terkait bidang tertentu. Kata benda, angka, kata ganti demonstratif seperti「\vocab{這}，\vocab{那}，\vocab{哪}，\vocab{每}，\vocab{各}，atau \vocab{別的}」dapat ditempatkan sebelum \vocab{方面}.\\
a. 我們員工對公司的福利都很滿意，可是在假期方面，還是有人覺得太少了。\\
b. 這次會議，在健康保險方面大家還有沒有什麼問題？
\end{enumerate}

\end{document}